% Plain-language ("simplified") version of the MIC paper, with the data
% tables (no figures) drawn from mic-paper-v9-ieee-tmi.tex.
% Compile: pdflatex mic-paper-simplified.tex (x2)
\documentclass[11pt]{article}

% ---- Page / layout ----
\usepackage[letterpaper,margin=1in]{geometry}
\usepackage{parskip}

% ---- Math / fonts ----
\usepackage{amsmath}
\usepackage{amssymb}

% ---- Graphics / tables ----
\usepackage{graphicx}
\usepackage{booktabs}
\usepackage{array}
\usepackage{multirow}

% ---- Utilities ----
\usepackage{xcolor}
\usepackage{url}
\usepackage[colorlinks=true,linkcolor=blue!70!black,urlcolor=blue!70!black]{hyperref}
\usepackage{microtype}

\title{A 16-Bit-Native Compression Pipeline for\\
       Lossless Medical Images}
\author{Kuldeep Singh}
\date{}

\begin{document}
\maketitle

\begin{quote}
\emph{A plain-language version of the MIC paper. It keeps the main ideas, the
key numbers, and the reasoning, but trades the formal equations and
implementation minutiae for explanations a software engineer can follow.
Data tables are reproduced from the full IEEE~TMI manuscript; figures are
omitted.}
\end{quote}

% ============================================================
\section*{Abstract}

Medical images have to be stored and transmitted without losing a single
pixel: a radiologist reading a scan must see exactly what the scanner
produced. These images are also ``deep'' --- each pixel is 10 to 16 bits,
not the 8 bits of an ordinary photo. That deep pixel format is where most
compression tools stumble, because almost all of them were built for
8-bit data.

This paper describes \textbf{MIC} (Medical Image Codec), a lossless codec that
works \emph{directly} on 16-bit data instead of chopping each pixel into two
bytes. MIC chains together three simple stages: a spatial predictor, a
16-bit run-length encoder, and a large-alphabet entropy coder called FSE
(Finite State Entropy, a fast table-driven form of arithmetic-style
coding). It also includes a ``multi-state'' decoder trick that keeps the CPU
busier and speeds up decompression.

On 21 real DICOM images across nine modalities, MIC:

\begin{itemize}
  \item compresses \textbf{5--22\% better} (14\% on average) than a strong
  general-purpose baseline (Delta + Zstandard);
  \item reaches an average compression ratio of \textbf{3.46}, basically tied
  with High-Throughput JPEG 2000 (HTJ2K, 3.45) and about 91\% of JPEG-LS
  (3.82);
  \item is the \textbf{fastest decoder} of the four codecs tested on 20 of 21
  images on ARM64 and on all 21 on AMD64;
  \item ships with a tiny \textbf{20\,KB pure-JavaScript decoder} so the format
  can be decoded right in a web browser.
\end{itemize}

The main caveat: this is a 21-image study. The results are encouraging but
need confirmation on larger, multi-hospital datasets.

% ============================================================
\section{Why Medical Images Need a Different Codec}

A modern scanner produces a lot of data. One digital breast tomosynthesis
view alone can be 60+ image slices and over 600\,MB; a full screening study
tops 2\,GB uncompressed. Hospitals store and move enormous volumes of this
data every day, so good lossless compression directly affects storage cost,
network load, and how fast images appear on a viewer.

DICOM, the medical imaging standard, already supports several lossless
formats --- JPEG 2000, HTJ2K, JPEG-LS, and a basic run-length scheme. Each
makes a different trade-off between how small the file gets, how fast it
decodes, and how complex it is to implement. JPEG 2000 and HTJ2K compress
well but rely on fairly heavy machinery. JPEG-LS usually gets the smallest
files. Plain run-length coding is simple but barely compresses.

\subsection{The core problem: everything assumes 8-bit data}

Here is the observation the whole paper is built on. The popular entropy
coders --- Huffman, arithmetic coding, the LZ77 family, and the FSE/ANS coder
inside Zstandard --- were all designed for \textbf{8-bit byte streams}, where
there are only 256 possible symbols. Medical images break that assumption: a
16-bit pixel has up to 65,536 possible values.

When an 8-bit coder meets 16-bit data, it has two bad options:

\begin{enumerate}
  \item \textbf{Split each pixel into a high byte and a low byte} and code them
  separately. This doubles the number of coding steps and, worse, throws
  away the relationship between the two bytes.
  \item \textbf{Cap the alphabet at 256 values} and store anything rarer as raw
  data, which wastes space.
\end{enumerate}

Why does splitting hurt? Consider a tiny prediction error like $-2$, $-1$, 0,
or $+1$. The low byte carries the real information, but the high byte is
almost always the same (just a sign-extension). A byte-by-byte coder still
spends bits encoding that high byte even though it's nearly fully
determined by the low one. On the test images, the information shared
between the two bytes ranges from 0.3 to 1.1 bits per pixel --- that is the
ceiling on how much a byte-splitting coder is wasting, and MIC's measured
5--22\% advantage sits comfortably under that ceiling.

\paragraph{The residual stream.} Throughout, a \emph{residual} is the
difference between a pixel and a prediction of it. Example: if the predictor
guesses 1020 and the true pixel is 1023, the residual is $+3$. Smooth regions
of an image produce lots of residuals near zero, and the same small values show
up again and again --- exactly the pattern run-length and entropy coding love.

\subsection{What MIC is}

MIC keeps the residuals at their native 16-bit width and combines:

\begin{enumerate}
  \item a \textbf{simple spatial predictor},
  \item a \textbf{16-bit run-length encoder}, and
  \item an \textbf{entropy coder (FSE) extended to handle large alphabets}.
\end{enumerate}

Plus a multi-state decoder that improves decompression speed by letting the
CPU work on several symbols at once.

\subsection{Contributions, briefly}

\begin{itemize}
  \item A \textbf{16-bit-native pipeline} (Delta + 16-bit RLE + extended FSE)
  that beats Delta + Zstandard on every one of the 21 images.
  \item An \textbf{entropy coder for large alphabets} --- up to 65,535 symbols,
  versus the 4,096-symbol limit in Zstandard's FSE. Tables shrink
  automatically to fit the data (as small as ${\sim}2$\,KB for 8-bit content).
  \item A \textbf{multi-state decoder} that runs 2/4/8 independent decode chains
  to use idle CPU capacity, with no cost to compression ratio.
  \item An \textbf{evaluation} against HTJ2K, JPEG-LS, and Delta + Zstandard.
  \item A \textbf{20\,KB JavaScript decoder} plus a WebAssembly build, for
  client-side decoding in browsers.
\end{itemize}

\paragraph{Scope.} This work covers lossless, single-frame, grayscale images
only. Color and multi-frame extensions, lossy modes, and progressive decoding
are out of scope.

% ============================================================
\section{Related Work, in One Page}

\textbf{The DICOM lossless codecs.} JPEG 2000 uses a wavelet transform plus a
sophisticated block coder (EBCOT) --- great ratios, slow decode. HTJ2K swaps
in a faster block coder and decodes much quicker. JPEG-LS uses a small
predictor (Median Edge Detector) and Golomb-Rice coding; it's a widely
deployed standard and a strong ratio baseline. The basic DICOM run-length
scheme barely compresses because it works on raw bytes with no prediction
step. Research codecs like CALIC and FLIF compress well on natural images
but aren't DICOM standards and lack browser decoders.

These are the codecs MIC is measured against. The goal isn't to dethrone
them, but to explore a \textbf{simpler, residual-domain design} aimed squarely
at high-bit-depth images, with strong speed and easy deployment.

\textbf{Entropy coding background.} Asymmetric Numeral Systems (ANS) is a
family of entropy coders that replaces arithmetic coding's interval math with
simple integer state updates. FSE is a fast, table-driven version of ANS
made popular by Zstandard. The catch: these implementations are tuned for
byte streams (Zstandard caps FSE at 4,096 symbols).

\textbf{Parallel ANS.} Prior work has shown ANS decoders go faster when you run
several independent streams at once (Giesen on interleaved coders; later
work scaling this to many streams and to GPUs). MIC doesn't claim to invent
ANS or interleaving --- it applies these known ideas inside a 16-bit-native
medical codec with a specific large-alphabet design.

Table~\ref{tab:positioning} places MIC against the prior work across five
dimensions relevant to medical image compression.

\begin{table}[htbp]
\centering
\caption{Feature comparison of MIC with related codecs and entropy coding
methods. ``Entropy coder sees 16-bit residual as one symbol'' is a property
of the entropy stage, not of the codec's overall bit-depth support: JPEG-LS
and HTJ2K both handle 16-bit medical images, but their entropy stages
(Golomb--Rice on context bins for JPEG-LS, bit-plane coding for HTJ2K) do
not treat the predicted residual as a single 16-bit symbol. Source-line
counts are order-of-magnitude context only; the reference libraries also
implement standard profiles outside MIC's current scope.}
\label{tab:positioning}
\footnotesize
\setlength{\tabcolsep}{3pt}
\begin{tabular}{p{3.2cm}ccccc}
\toprule
Feature & JPEG-LS & HTJ2K & Zstd & rANS & MIC \\
\midrule
Lossless medical & Yes & Yes & General & General & Yes \\
Entropy coder sees 16-bit residual as one symbol & No & No & No & Configurable & Yes \\
Multi-state ILP & No & No & No & Yes & Yes \\
Browser decoder & Via WASM port & Via WASM port & Partial & No & Purpose-built JS \\
Source lines & 5k & 20k & N/A & N/A & 1.1k \\
\bottomrule
\end{tabular}
\end{table}

% ============================================================
\section{How MIC Works}

The pipeline is three stages, each simple enough to decode in one
straight-line pass with very few branches:

\begin{verbatim}
Raw 16-bit pixels
   -> Spatial prediction  (subtract a guess based on neighbors)
   -> 16-bit run-length encoding  (collapse repeated values)
   -> FSE entropy coding  (squeeze out the remaining redundancy)
   -> Compressed bytes
\end{verbatim}

\subsection{Spatial prediction}

For each pixel, MIC predicts its value as the \textbf{average of the pixel
above and the pixel to the left}, then stores only the difference (the
residual). Pixels on the top edge or left edge use whatever single neighbor is
available; the very first pixel is stored as-is. Division rounds down, and
the decoder repeats the exact same arithmetic so reconstruction is exact.

This predictor was chosen because it's cheap, has almost no branches on the
decode side, and works well on medical images.

\paragraph{Other predictors we tried.} We compared four predictors (left-only,
average, Paeth from PNG, and MED from JPEG-LS) through the full pipeline.
MED gave the best ratio (about 1.6\% better than average) but decoded
1.5--2$\times$ slower because of its three-way branch and a diagonal dependency
that blocks the fast branch-free loop. The average predictor won on
consistency and speed, so it's the default. Table~\ref{tab:predictor-summary}
summarizes the comparison.

\begin{table}[htbp]
\centering
\caption{Predictor ablation summary: geometric means and win counts across all
21 images.}
\label{tab:predictor-summary}
\begin{tabular}{lrrr}
\toprule
Predictor & Geo.\ Mean & Wins & Avg$\to$X \\
\midrule
Left-only & 3.38$\times$ & 3/21 & $-$2.3\% \\
Average (MIC) & 3.46$\times$ & 13/21 & baseline \\
Paeth & 3.48$\times$ & 2/21 & $+$0.5\% \\
MED (JPEG-LS) & 3.52$\times$ & 16/21 & $+$1.6\% \\
\bottomrule
\end{tabular}
\end{table}

\subsection{Handling big jumps: overflow coding}

Most residuals are small, but occasionally a pixel jumps far from its
prediction. Rather than widen the whole stream to 32 bits to handle rare big
values, MIC reserves one special code as a \textbf{delimiter}:

\begin{itemize}
  \item Small, in-range residuals are stored directly as 16-bit codes (with
  zero mapped to a fixed midpoint value).
  \item A large, out-of-range residual is stored as the delimiter followed by
  the raw pixel value.
\end{itemize}

On decode, the reader checks each value: if it's the delimiter, the next
value is a raw pixel; otherwise it's a normal residual. The ranges are
designed not to overlap, so there's never any ambiguity. In normal 8--16-bit
clinical data the overflow path is rare; it exists to guarantee correctness,
not for the common case. Table~\ref{tab:overflow-map} shows the mapping for a
representative bit depth.

\begin{table}[htbp]
\centering
\caption{Overflow coding example for $d=12$ bits
(\texttt{deltaThreshold}$=2047$, delimiter $D=4095$).}
\label{tab:overflow-map}
\begin{tabular}{lll}
\toprule
Residual $r$ & Emitted code(s) & Notes \\
\midrule
$-2$ & $2045$ & in-range \\
$-1$ & $2046$ & in-range \\
$0$  & $2047$ & in-range (= threshold) \\
$+1$ & $2048$ & in-range \\
$+2$ & $2049$ & in-range \\
$\pm 2047$ & $4095$, raw pixel & overflow \\
\bottomrule
\end{tabular}
\end{table}

\subsection{16-bit run-length encoding (RLE)}

The residual stream is turned into two kinds of runs:

\begin{itemize}
  \item \textbf{Same runs:} a count plus one repeated value (used when a value
  appears at least three times in a row).
  \item \textbf{Literal runs:} a count followed by that many values copied
  verbatim (the values need not be distinct).
\end{itemize}

A single header number tells the two apart using a midpoint trick: small
header values mean ``same run of this length,'' large ones mean ``literal run.''
The longest single run is ${\sim}32{,}767$ symbols; longer stretches are split
across several headers automatically.

\textbf{Why 16-bit RLE matters.} Suppose 1,000 identical 16-bit residuals
appear in a row. MIC encodes that as just two numbers: a count and the value.
A \emph{byte}-oriented RLE sees the high and low bytes interleaved, so no
single byte repeats 1,000 times --- it's forced into two interleaved runs or
falls back to literals. Working at 16 bits captures the repetition that
byte-splitting hides.

\textbf{Fast decode.} Same runs (often 80\% of all runs on smooth images) are
fast-pathed: the decoder just hands back the cached value and decrements a
counter, without touching the compressed data at all.

\subsection{The entropy coder (FSE) for large alphabets}

\textbf{FSE in one paragraph.} FSE keeps a single integer ``state.'' To decode
one symbol, it looks up the current state in a table; the entry tells it which
symbol to output, how many bits to read next, and how to compute the next
state. So each symbol costs \textbf{one table lookup, one bit read, and one
addition} --- no multiplies or divides. (Encoding runs the symbols in reverse
and builds the bitstream backwards; decoding reads forwards. That reversal
is just how ANS works.)

MIC extends FSE to handle up to \textbf{65,535 symbols} instead of Zstandard's
4,096, because high-bit-depth residuals genuinely need a big alphabet. It
also \textbf{sizes its tables to the data}: an 8-bit image stored in 16-bit
containers (common for MR) shrinks the working tables from ${\sim}512$\,KB down
to about 2\,KB, which keeps everything in fast cache.

\textbf{Picking the table size automatically.} FSE has a ``tableLog'' knob that
sets how big and precise the probability table is. MIC picks it based on how
many distinct values it sees and how much data backs each value: it uses a
larger table only when there are many active symbols \emph{and} enough data to
estimate their probabilities reliably (for example, mammography images after
prediction). On the test set this adaptive choice correctly turned on the
larger table for exactly the nine images that benefit (1--10\% better ratio),
and decode speed is unaffected by the choice.
Table~\ref{tab:tablelog-ablation} shows the ablation across selected images.

\begin{table}[htbp]
\centering
\caption{TableLog ablation: forced TL=11/12/13 vs.\ adaptive selection
(selected images, ARM64 reference platform).}
\label{tab:tablelog-ablation}
\begin{tabular}{lrrrr}
\toprule
Image & TL=11 & TL=12 & TL=13 & Adaptive \\
\midrule
CR  & 3.47$\times$ & 3.63$\times$ & 3.69$\times$ & 3.69$\times$ \\
MG1 & 8.00$\times$ & 8.57$\times$ & 8.79$\times$ & 8.79$\times$ \\
MG2 & 7.98$\times$ & 8.55$\times$ & 8.77$\times$ & 8.77$\times$ \\
MR2 & 3.14$\times$ & 3.24$\times$ & 3.28$\times$ & 3.28$\times$ \\
RG2 & 3.85$\times$ & 4.12$\times$ & 4.23$\times$ & 4.23$\times$ \\
RG3 & 5.64$\times$ & 5.95$\times$ & 6.08$\times$ & 6.08$\times$ \\
\bottomrule
\end{tabular}
\end{table}

% ============================================================
\section{Going Faster: Multi-State Decoding}

\subsection{The bottleneck}

A normal single-state ANS decoder has a problem: each symbol depends on the
one before it. To decode symbol $N$ you need the state produced by symbol
$N-1$, which means the work can't start early. With a table lookup taking
roughly 4--5 CPU cycles, the loop is stuck at about one symbol every few
cycles --- even though modern CPUs have 4--6 execution ports sitting mostly
idle. A single decode chain uses just one of them, leaving ${\sim}75$--83\% of
the CPU's capacity unused.

\subsection{The fix}

MIC runs \textbf{several independent decoders at once} on interleaved
positions. With 8 chains, chain 0 handles positions $0, 8, 16, \ldots$; chain 1
handles $1, 9, 17, \ldots$; and so on. The chains share nothing during decoding
--- each has its own state and its own stream of lookups --- so the CPU's
out-of-order engine can keep all of them in flight simultaneously.

On the encoder side, the symbols are split into $N$ interleaved groups, each
compressed independently, and the resulting streams are concatenated. They
all share one decode table, and each chain's starting state is recorded in
the header. The decoder just writes each chain's output back to its
positions.

\subsection{How much it helps}

In theory, $N$ chains could be up to $N\times$ faster. In practice the speedup
is smaller because the chains share some bookkeeping (refilling the bit
reader), the unrolled loop puts pressure on the instruction cache, and reading
the compressed data has its own bandwidth limit. Measured on the FSE stage
alone:

\begin{itemize}
  \item \textbf{4 states:} $+$46\% on ARM64, $+$8\% on AMD64 (vs.\ single
  state)
  \item \textbf{8 states:} a further $+$7\% on ARM64, $+$2\% on AMD64
\end{itemize}

with \textbf{no change to the compressed file}. ARM64 keeps improving as chains
double; AMD64 flattens out earlier because its core already extracts a lot of
parallelism from a single chain on its own. Table~\ref{tab:latency} shows the
simple latency model against the measured ARM64 speedup.

\begin{table}[htbp]
\centering
\caption{Multi-state decoder latency model and empirical speedup.}
\label{tab:latency}
\begin{tabular}{clrr}
\toprule
States ($k$) & Amortized latency & Theoretical & Empirical (ARM64) \\
\midrule
1 & $L$ cy/sym & 1.0$\times$ & 1.0$\times$ \\
2 & $L/2$ cy/sym & 2.0$\times$ & 1.3$\times$ \\
4 & $L/4$ cy/sym & 4.0$\times$ & 1.5$\times$ \\
8 & $L/8$ cy/sym & 8.0$\times$ & 1.6$\times$ \\
\bottomrule
\end{tabular}
\end{table}

The decoder signals its mode with a short magic header (a reserved byte
value that can't be a valid single-state setting), so single-, two-, four-,
and eight-state files all coexist without confusion. On AMD64 and ARM64 the
inner loop uses each architecture's native variable-shift instructions; other
platforms get a pure-Go fallback. Table~\ref{tab:format} summarises the
on-disk layout of a single-frame stream.

\begin{table}[htbp]
\centering
\caption{Single-frame MIC bitstream layout. All multi-byte fields are
little-endian. \emph{States} $k$ is the number of interleaved FSE chains
($k \in \{1,2,4,8\}$).}
\label{tab:format}
\footnotesize
\begin{tabular}{p{2.6cm}rp{6.0cm}}
\toprule
Field & Size & Meaning \\
\midrule
Magic / mode byte & 1 B & Single-state: encodes \texttt{tableLog} (value 5--17). Multi-state: \texttt{0xFF}. \\
Multi-state tag & 1 B & For multi-state streams, one of \texttt{0x02}, \texttt{0x04}, \texttt{0x84} indicating $k=2$, 4, or 8. \\
Symbol count & 4 B & Number of decoded RLE symbols. \\
TableLog & 1 B & Decode-table size: $2^{\textit{tableLog}}$ entries (multi-state path only; the single-state path overloads byte~0). \\
Normalized counts & variable & FSE normalised symbol frequencies, serialised with the standard variable-length encoding. \\
Per-chain initial states & $k \times 2$ B & Starting state for each FSE chain. \\
Per-chain stream lengths & $k \times 4$ B & Compressed-stream length per chain (omitted when $k=1$). \\
Per-chain compressed streams & variable & Concatenated chain payloads. \\
\bottomrule
\end{tabular}
\end{table}

% ============================================================
\section{How We Measured It}

\textbf{Dataset.} 21 de-identified DICOM images across nine modalities (CT, MR,
CR, X-ray, mammography, tomosynthesis, and others), from public NEMA sample
sets and a public breast-tomosynthesis case. It's a varied set but small
relative to a real clinical archive --- hence ``on this 21-image dataset''
rather than sweeping claims. Table~\ref{tab:dataset} lists the images.

\begin{table}[htbp]
\centering
\caption{Test dataset: 21 clinical DICOM images spanning nine modalities.}
\label{tab:dataset}
\footnotesize
\begin{tabular}{llrr}
\toprule
Image & Modality & Dimensions & Raw (MB) \\
\midrule
MR   & Brain MRI      & $256{\times}256$   & 0.13 \\
CT   & CT scan        & $512{\times}512$   & 0.50 \\
CR   & Comp.\ radiography & $2140{\times}1760$ & 7.18 \\
XR   & X-ray          & $2048{\times}2577$ & 10.1 \\
MG1  & Mammography    & $2457{\times}1996$ & 9.35 \\
MG2  & Mammography    & $2457{\times}1996$ & 9.35 \\
MG3  & Mammography    & $4774{\times}3064$ & 27.3 \\
MG4  & Mammography    & $4096{\times}3328$ & 26.0 \\
CT1  & CT scan        & $512{\times}512$   & 0.50 \\
CT2  & CT scan        & $512{\times}512$   & 0.50 \\
MG-N & Mammography    & $3064{\times}4664$ & 27.3 \\
MR1  & Brain MRI      & $512{\times}512$   & 0.50 \\
MR2  & Brain MRI      & $1024{\times}1024$ & 2.00 \\
MR3  & Brain MRI      & $512{\times}512$   & 0.50 \\
MR4  & Brain MRI      & $512{\times}512$   & 0.50 \\
NM1  & Nuclear med.   & $256{\times}1024$  & 0.50 \\
RG1  & Radiography    & $1841{\times}1955$ & 6.86 \\
RG2  & Radiography    & $1760{\times}2140$ & 7.18 \\
RG3  & Radiography    & $1760{\times}1760$ & 5.91 \\
SC1  & Sec.~capture   & $2048{\times}2487$ & 9.71 \\
XA1  & Fluoroscopy    & $1024{\times}1024$ & 2.00 \\
\bottomrule
\end{tabular}
\end{table}

\textbf{Baselines.} HTJ2K (via OpenJPH), JPEG-LS (via CharLS), and Delta +
Zstandard at its maximum level. All were called in-process (no file or
subprocess overhead) for a fair speed comparison.

\textbf{Metrics.} Compression ratio (uncompressed $\div$ compressed), and
decode and encode throughput in MB/s of uncompressed pixels.

\textbf{Platforms.} Two reference machines: an ARM64 (Apple M4 Pro) and an
AMD64 (AWS Intel Xeon 6). Go code used the standard compiler; C components were
built with \texttt{-O3}. Each benchmark ran the full image 10 times;
run-to-run variance stayed under 2\% on ARM64 and 5\% on AMD64.

The ``MIC'' column in the result tables is the fastest available C decoder on
each platform --- the eight-state decoder, with AVX-512 acceleration of the
run-length and prediction-reversal stages on AMD64. Both read identical
compressed bytes; only the vector width differs.

\textbf{Browser decoder.} A 20\,KB pure-JavaScript ES module (zero npm
dependencies) plus a 2.5\,MB Go WebAssembly build. The JS decoder
auto-detects 1/2/4/8-state streams and was benchmarked under Node.js on the
ARM64 machine.

% ============================================================
\section{Results}

\subsection{Compression ratio}

JPEG-LS gets the smallest files on every image --- its context-adaptive
predictor is hard to beat on pure ratio. But MIC isn't optimizing ratio
alone; it's after the best balance of ratio, speed, and simplicity. Across
the dataset MIC averages a \textbf{3.46} ratio: about 91\% of JPEG-LS (3.82)
and essentially tied with HTJ2K (3.45), while decoding faster than both.
Mammography compresses best of all (up to 8.79$\times$) because its smooth
tissue produces long near-zero runs. Table~\ref{tab:ratios} gives the
per-image ratios for all variants.

\begin{table}[htbp]
\centering
\caption{Lossless compression ratios: all codec variants, 21 images. Bold =
best ratio per row.}
\label{tab:ratios}
\footnotesize
\setlength{\tabcolsep}{4pt}
\resizebox{\textwidth}{!}{%
\begin{tabular}{lrrrrrrrr}
\toprule
Image & Raw (MB) & $\Delta$+Zstd-19 & MIC & Wavelet & PICS-4 & PICS-8 & HTJ2K & JPEG-LS \\
\midrule
MR   & 0.13 & 1.95$\times$ & 2.35$\times$ & 2.38$\times$ & 2.28$\times$ & 2.21$\times$ & 2.38$\times$ & \textbf{2.52$\times$} \\
CT   & 0.50 & 2.05$\times$ & 2.24$\times$ & 1.67$\times$ & 2.15$\times$ & 1.96$\times$ & 1.77$\times$ & \textbf{2.68$\times$} \\
CR   & 7.18 & 3.24$\times$ & 3.69$\times$ & 3.81$\times$ & 3.70$\times$ & 3.71$\times$ & 3.77$\times$ & \textbf{3.96$\times$} \\
XR   & 10.1 & 1.43$\times$ & 1.74$\times$ & \textbf{1.76$\times$} & 1.75$\times$ & \textbf{1.76$\times$} & 1.67$\times$ & \textbf{1.76$\times$} \\
MG1  & 9.35 & 7.20$\times$ & 8.79$\times$ & 8.67$\times$ & 8.84$\times$ & 8.87$\times$ & 8.25$\times$ & \textbf{8.91$\times$} \\
MG2  & 9.35 & 7.21$\times$ & 8.77$\times$ & 8.65$\times$ & 8.83$\times$ & 8.85$\times$ & 8.24$\times$ & \textbf{8.90$\times$} \\
MG3  & 27.3 & 2.04$\times$ & 2.24$\times$ & 2.32$\times$ & 2.31$\times$ & 2.34$\times$ & 2.22$\times$ & \textbf{2.38$\times$} \\
MG4  & 26.0 & 3.10$\times$ & 3.47$\times$ & 3.59$\times$ & 3.59$\times$ & 3.62$\times$ & 3.51$\times$ & \textbf{3.71$\times$} \\
CT1  & 0.50 & 2.47$\times$ & 2.79$\times$ & 2.49$\times$ & 2.54$\times$ & 2.29$\times$ & 2.70$\times$ & \textbf{3.19$\times$} \\
CT2  & 0.50 & 3.24$\times$ & 3.49$\times$ & 2.87$\times$ & 3.11$\times$ & 2.72$\times$ & 3.29$\times$ & \textbf{4.54$\times$} \\
MG-N & 27.3 & 2.04$\times$ & 2.24$\times$ & 2.32$\times$ & 2.31$\times$ & 2.34$\times$ & 2.23$\times$ & \textbf{2.38$\times$} \\
MR1  & 0.50 & 1.79$\times$ & 2.09$\times$ & 2.14$\times$ & 2.10$\times$ & 2.08$\times$ & 2.13$\times$ & \textbf{2.30$\times$} \\
MR2  & 2.00 & 2.77$\times$ & 3.28$\times$ & 3.34$\times$ & 3.31$\times$ & 3.31$\times$ & 3.35$\times$ & \textbf{3.52$\times$} \\
MR3  & 0.50 & 3.47$\times$ & 3.93$\times$ & 4.09$\times$ & 3.89$\times$ & 3.84$\times$ & 4.33$\times$ & \textbf{4.51$\times$} \\
MR4  & 0.50 & 3.58$\times$ & 4.12$\times$ & 4.18$\times$ & 4.09$\times$ & 4.03$\times$ & 4.21$\times$ & \textbf{4.49$\times$} \\
NM1  & 0.50 & 4.88$\times$ & 5.15$\times$ & 5.02$\times$ & 5.26$\times$ & 5.28$\times$ & 5.76$\times$ & \textbf{6.28$\times$} \\
RG1  & 6.86 & 1.45$\times$ & 1.70$\times$ & 1.70$\times$ & 1.70$\times$ & 1.69$\times$ & 1.63$\times$ & \textbf{1.72$\times$} \\
RG2  & 7.18 & 3.68$\times$ & 4.23$\times$ & 4.32$\times$ & 4.28$\times$ & 4.30$\times$ & 4.32$\times$ & \textbf{4.51$\times$} \\
RG3  & 5.91 & 5.46$\times$ & 6.08$\times$ & 6.82$\times$ & 6.11$\times$ & 6.12$\times$ & 6.99$\times$ & \textbf{7.31$\times$} \\
SC1  & 9.71 & 3.46$\times$ & 3.71$\times$ & 3.70$\times$ & 3.73$\times$ & 3.74$\times$ & 3.85$\times$ & \textbf{4.73$\times$} \\
XA1  & 2.00 & 4.37$\times$ & 5.01$\times$ & 4.94$\times$ & 5.04$\times$ & 5.03$\times$ & 4.88$\times$ & \textbf{5.39$\times$} \\
\midrule
\textbf{Geomean} & --- & \textbf{3.03$\times$} & \textbf{3.46$\times$} & \textbf{3.41$\times$} & \textbf{3.44$\times$} & \textbf{3.38$\times$} & \textbf{3.45$\times$} & \textbf{3.82$\times$} \\
\bottomrule
\end{tabular}}
\end{table}

\subsection{Encoding speed}

MIC encodes faster than the others on most images, on both platforms,
because its encoder is a single pass (predict, run-length, table-driven
entropy code) while HTJ2K and JPEG-LS do more work per pixel.

\begin{itemize}
  \item \textbf{ARM64:} MIC wins on all 21 images --- roughly 1.9--3.2$\times$
  HTJ2K, 3--7$\times$ JPEG-LS, and 35--240$\times$ Delta + Zstandard (whose
  max-level LZ77 search is what makes it so slow).
  \item \textbf{AMD64:} MIC wins on 18 of 21; HTJ2K narrowly takes the two
  biggest mammography slabs and one radiography image. The gap is smaller here
  because the AMD64 baselines are themselves heavily vectorized.
\end{itemize}

Table~\ref{tab:enc} gives the per-image encoding throughput.

\begin{table}[htbp]
\centering
\caption{Single-threaded encoding throughput (MB/s). MIC is the
eight-state C encoder. Bold marks the fastest codec per row on each
platform.}
\label{tab:enc}
\footnotesize
\setlength{\tabcolsep}{4pt}
\resizebox{\textwidth}{!}{%
\begin{tabular}{l|rrrr|rrrr}
\toprule
 & \multicolumn{4}{c|}{ARM64} & \multicolumn{4}{c}{AMD64} \\
\cmidrule(lr){2-5}\cmidrule(lr){6-9}
Image & MIC & $\Delta$+Zstd & HTJ2K & JPEG-LS & MIC & $\Delta$+Zstd & HTJ2K & JPEG-LS \\
\midrule
MR   & \textbf{504}   &  8 & 273 &  97 & \textbf{338} &  5 & 229          &  62 \\
CT   & \textbf{529}   &  8 & 200 & 139 & \textbf{292} &  5 & 206          &  97 \\
CR   & \textbf{569}   &  3 & 204 & 115 & \textbf{322} &  2 & 284          &  75 \\
XR   & \textbf{690}   &  3 & 248 & 125 & \textbf{405} &  3 & 250          &  84 \\
MG1  & \textbf{1{,}168} & 14 & 456 & 315 & 546          &  8 & \textbf{647} & 209 \\
MG2  & \textbf{1{,}172} & 14 & 454 & 315 & 541          &  8 & \textbf{646} & 209 \\
MG3  & \textbf{706}   &  3 & 234 & 143 & \textbf{332} &  2 & 251          & 104 \\
MG4  & \textbf{993}   &  7 & 373 & 215 & \textbf{431} &  6 & 415          & 150 \\
CT1  & \textbf{576}   & 11 & 238 & 178 & \textbf{329} &  7 & 277          & 125 \\
CT2  & \textbf{561}   &  8 & 207 & 172 & \textbf{286} &  5 & 266          & 125 \\
MG-N & \textbf{718}   &  2 & 235 & 143 & \textbf{339} &  2 & 243          & 104 \\
MR1  & \textbf{677}   & 11 & 215 & 120 & \textbf{344} &  6 & 244          &  86 \\
MR2  & \textbf{737}   &  9 & 231 & 131 & \textbf{397} &  4 & 300          &  84 \\
MR3  & \textbf{901}   & 22 & 299 & 187 & \textbf{428} & 12 & 353          & 123 \\
MR4  & \textbf{711}   &  8 & 228 & 175 & \textbf{358} &  6 & 318          & 141 \\
NM1  & \textbf{663}   &  7 & 247 & 170 & \textbf{336} &  5 & 312          & 116 \\
RG1  & \textbf{659}   &  5 & 246 &  95 & \textbf{392} &  4 & 250          &  65 \\
RG2  & \textbf{811}   &  6 & 271 & 151 & \textbf{416} &  3 & 356          &  99 \\
RG3  & \textbf{762}   &  6 & 285 & 189 & 397          &  4 & \textbf{407} & 126 \\
SC1  & \textbf{765}   &  6 & 256 & 223 & \textbf{407} &  4 & 309          & 168 \\
XA1  & \textbf{693}   &  7 & 232 & 152 & \textbf{358} &  4 & 310          &  97 \\
\midrule
\textbf{Geomean} & \textbf{740} & \textbf{7} & \textbf{273} & \textbf{161} & \textbf{376} & \textbf{4} & \textbf{311} & \textbf{110} \\
\bottomrule
\end{tabular}}
\end{table}

\subsection{Decoding speed}

This is MIC's strongest result. Its decoder is dominated by that tight FSE
loop --- one lookup, one bit read, one add --- with no hard-to-predict
branches. HTJ2K and JPEG-LS both have data-dependent branches that cause CPU
mispredictions. Zstandard is the closest relative (same predictor, also uses
FSE) but works on bytes, so it pays the 2$\times$ byte-splitting cost on deep
data.

\begin{itemize}
  \item \textbf{ARM64:} MIC is fastest on \textbf{20 of 21} images (Zstandard
  edges it out only on CT, by 1.3\%). Roughly 1.5--2.6$\times$ HTJ2K,
  1.1--1.7$\times$ Zstandard, and 2.0--4.5$\times$ JPEG-LS.
  \item \textbf{AMD64:} MIC is fastest on \textbf{all 21} images, from a hair
  ahead on mammography up to 1.62$\times$ on X-ray. Switching from four to
  eight states (plus the AVX-512 widening of the post-entropy stages) was what
  flipped the few remaining close calls into MIC wins.
\end{itemize}

So MIC delivers both \textbf{better ratio than Zstandard on every image} and
\textbf{faster decode on nearly every image} --- the signature of the
16-bit-native design. JPEG-LS keeps about a 10\% ratio edge, so the real
choice between them is decode speed versus storage cost. Per-image numbers are
in Table~\ref{tab:decomp}, and the entropy stage in isolation is in
Table~\ref{tab:fse-combined}.

\begin{table}[htbp]
\centering
\caption{Single-threaded decompression throughput (MB/s). MIC is
the eight-state C decoder (AVX-512 acceleration of the RLE and delta
inverse stages on AMD64). Bold marks the fastest codec per row on
each platform.}
\label{tab:decomp}
\footnotesize
\setlength{\tabcolsep}{4pt}
\resizebox{\textwidth}{!}{%
\begin{tabular}{l|rrrr|rrrr}
\toprule
 & \multicolumn{4}{c|}{ARM64} & \multicolumn{4}{c}{AMD64} \\
\cmidrule(lr){2-5}\cmidrule(lr){6-9}
Image & MIC & $\Delta$+Zstd & HTJ2K & JPEG-LS & MIC & $\Delta$+Zstd & HTJ2K & JPEG-LS \\
\midrule
MR   & \textbf{564} & 331          & 301 & 146 & \textbf{552} & 357 & 439 &  91 \\
CT   & 466          & \textbf{472} & 356 & 187 & \textbf{453} & 384 & 329 & 111 \\
CR   & \textbf{692} & 564          & 367 & 206 & \textbf{566} & 475 & 439 & 123 \\
XR   & \textbf{655} & 572          & 366 & 150 & \textbf{617} & 508 & 386 &  88 \\
MG1  & \textbf{829} & 691          & 681 & 547 & \textbf{913} & 599 & 907 & 332 \\
MG2  & \textbf{873} & 684          & 686 & 557 & \textbf{920} & 599 & 896 & 333 \\
MG3  & \textbf{678} & 473          & 362 & 201 & \textbf{556} & 358 & 390 & 117 \\
MG4  & \textbf{815} & 581          & 533 & 259 & \textbf{676} & 466 & 611 & 157 \\
CT1  & \textbf{571} & 502          & 373 & 238 & \textbf{525} & 425 & 390 & 143 \\
CT2  & \textbf{648} & 575          & 372 & 222 & \textbf{512} & 456 & 383 & 133 \\
MG-N & \textbf{657} & 469          & 360 & 201 & \textbf{545} & 355 & 391 & 117 \\
MR1  & \textbf{681} & 437          & 354 & 154 & \textbf{603} & 384 & 397 &  94 \\
MR2  & \textbf{747} & 505          & 390 & 224 & \textbf{644} & 464 & 491 & 132 \\
MR3  & \textbf{859} & 507          & 453 & 329 & \textbf{773} & 485 & 543 & 192 \\
MR4  & \textbf{775} & 523          & 304 & 230 & \textbf{632} & 480 & 446 & 153 \\
NM1  & \textbf{866} & 595          & 404 & 277 & \textbf{598} & 486 & 467 & 172 \\
RG1  & \textbf{561} & 413          & 354 & 126 & \textbf{550} & 347 & 363 &  80 \\
RG2  & \textbf{755} & 574          & 408 & 255 & \textbf{730} & 497 & 535 & 154 \\
RG3  & \textbf{740} & 655          & 469 & 320 & \textbf{723} & 552 & 628 & 196 \\
SC1  & \textbf{696} & 587          & 384 & 280 & \textbf{690} & 484 & 456 & 183 \\
XA1  & \textbf{759} & 611          & 388 & 278 & \textbf{688} & 534 & 477 & 165 \\
\midrule
\textbf{Geomean} & \textbf{733} & \textbf{557} & \textbf{386} & \textbf{227} & \textbf{602} & \textbf{501} & \textbf{475} & \textbf{144} \\
\bottomrule
\end{tabular}}
\end{table}

\begin{table}[htbp]
\centering
\caption{FSE decompression throughput (MB/s) for representative modalities,
single-thread pure-Go decoder. Increasing the number of independent state
chains exposes more instruction-level parallelism to the out-of-order
engine.}
\label{tab:fse-combined}
\footnotesize
\setlength{\tabcolsep}{6pt}
\begin{tabular}{lrrrrrr}
\toprule
 & \multicolumn{3}{c}{ARM64 (MB/s)} & \multicolumn{3}{c}{AMD64 (MB/s)} \\
\cmidrule(lr){2-4}\cmidrule(lr){5-7}
Image & 1-st & 4-st & 8-st & 1-st & 4-st & 8-st \\
\midrule
MR   &   612 &   953 &   765 &   636 &   899 &   759 \\
CT   &   776 & 1,872 & 1,729 &   974 &   986 &   953 \\
CR   & 3,051 & 5,749 & 5,521 & 1,426 & 1,581 & 1,638 \\
XR   & 3,657 & 6,072 & 5,807 & 1,939 & 1,890 & 2,248 \\
MG1  & 3,604 & 4,557 & 5,887 & 2,075 & 1,521 & 1,753 \\
MG-N & 3,588 & 6,378 & 6,307 & 1,714 & 2,294 & 2,473 \\
NM1  & 1,798 & 2,535 & 2,740 & 1,091 & 1,274 & 1,522 \\
RG1  & 1,786 & 3,492 & 5,412 & 1,198 & 1,717 & 2,104 \\
\midrule
\textbf{Geomean (21)} & \textbf{2,349} & \textbf{3,422} & \textbf{3,667} &
\textbf{1,464} & \textbf{1,575} & \textbf{1,614} \\
\bottomrule
\end{tabular}
\end{table}

\subsection{Other MIC variants (for reference)}

The main tables use one MIC column, but the codebase has several variants
that exist for different deployment needs rather than to compete for the
fastest number:

\begin{itemize}
  \item \textbf{Pure Go:} the whole thing in Go with no C dependency, at
  roughly 40--60\% of the C decoder's speed. Useful where you can't link C
  (mobile, WebAssembly, locked-down runtimes).
  \item \textbf{1/2/4-state decoders:} the same C decoder with fewer chains;
  they exist as baselines for the multi-state comparison.
  \item \textbf{Wavelet-based MIC:} swaps the predictor for a reversible 5/3
  wavelet transform feeding the same back end. Competitive on a couple of
  images but generally slower on this corpus, since smooth medical images
  don't reward the wavelet's extra cost. Kept as an option, not in the main
  tables.
\end{itemize}

The eight-state C decoder ends up being both the simplest and the fastest
configuration on this dataset.

\subsection{Multi-threaded decoding}

MIC can also split an image into horizontal strips and decode them on
separate threads. This scales nearly linearly up to four threads on both
platforms and keeps scaling to eight threads on images above ${\sim}5$\,MB,
hitting \textbf{4--4.6\,GB/s} on the largest mammography images. Small images
don't benefit --- below about half a megabyte, thread setup costs more than it
saves. We report this separately because the other codecs run single-threaded
in our pipeline, so a head-to-head wouldn't be fair.
Table~\ref{tab:pics} gives the per-image strip-parallel numbers.

\begin{table}[htbp]
\centering
\caption{Multi-threaded decompression throughput (MB/s) for MIC's
strip-parallel mode. The ``MIC'' column reproduces the single-threaded
baseline from Table~\ref{tab:decomp} exactly; $N$-strip columns decode the
same compressed bytes in $N$ parallel threads. The MR image is too small
for 8-way parallelism on both platforms ($^\dagger$).}
\label{tab:pics}
\footnotesize
\setlength{\tabcolsep}{4pt}
\resizebox{\textwidth}{!}{%
\begin{tabular}{l|rrrr|rrrr}
\toprule
 & \multicolumn{4}{c|}{ARM64} & \multicolumn{4}{c}{AMD64} \\
\cmidrule(lr){2-5}\cmidrule(lr){6-9}
Image & MIC & 2-strip & 4-strip & 8-strip & MIC & 2-strip & 4-strip & 8-strip \\
\midrule
MR   & 564   & 663   & 894     & 702\,$^\dagger$   & 552   & 329     & 367     & 294\,$^\dagger$ \\
CT   & 466   & 624   & 1{,}157 & 1{,}446           & 453   & 398     & 633     & 714 \\
CR   & 692   & 1{,}052 & 1{,}976 & 3{,}605         & 566   & 758     & 1{,}320 & 2{,}169 \\
XR   & 655   & 1{,}087 & 2{,}024 & 3{,}794         & 617   & 793     & 1{,}515 & 2{,}579 \\
MG1  & 829   & 1{,}529 & 2{,}529 & 4{,}443         & 913   & 1{,}317 & 1{,}951 & 2{,}890 \\
MG2  & 873   & 1{,}417 & 2{,}506 & 4{,}585         & 920   & 1{,}356 & 1{,}971 & 3{,}026 \\
MG3  & 678   & 1{,}112 & 1{,}911 & 3{,}820         & 556   & 838     & 1{,}449 & 2{,}767 \\
MG4  & 815   & 1{,}326 & 2{,}345 & 4{,}348         & 676   & 1{,}207 & 1{,}897 & 3{,}126 \\
CT1  & 571   & 819   & 1{,}247 & 1{,}525           & 525   & 458     & 703     & 725 \\
CT2  & 648   & 864   & 1{,}260 & 1{,}473           & 512   & 481     & 734     & 719 \\
MG-N & 657   & 1{,}147 & 2{,}027 & 3{,}930         & 545   & 913     & 1{,}537 & 2{,}819 \\
MR1  & 681   & 977   & 1{,}460 & 1{,}851           & 603   & 504     & 761     & 764 \\
MR2  & 747   & 1{,}137 & 1{,}911 & 3{,}103         & 644   & 778     & 1{,}171 & 1{,}534 \\
MR3  & 859   & 1{,}205 & 1{,}693 & 2{,}089         & 773   & 591     & 788     & 835 \\
MR4  & 775   & 997   & 1{,}469 & 1{,}757           & 632   & 540     & 796     & 804 \\
NM1  & 866   & 1{,}147 & 1{,}682 & 2{,}133         & 598   & 523     & 824     & 762 \\
RG1  & 561   & 677   & 1{,}203 & 2{,}150           & 550   & 603     & 1{,}101 & 1{,}890 \\
RG2  & 755   & 1{,}229 & 2{,}156 & 3{,}816         & 730   & 916     & 1{,}473 & 2{,}268 \\
RG3  & 740   & 1{,}335 & 2{,}343 & 4{,}168         & 723   & 971     & 1{,}618 & 2{,}730 \\
SC1  & 696   & 1{,}311 & 2{,}206 & 3{,}971         & 690   & 1{,}024 & 1{,}673 & 2{,}541 \\
XA1  & 759   & 1{,}214 & 1{,}881 & 3{,}191         & 688   & 792     & 1{,}246 & 1{,}748 \\
\bottomrule
\end{tabular}}
\end{table}

Practical rule of thumb: decode many small images? Use single-threaded MIC
per request (already fast). Decode one big image on a multi-core box? Use
strip-parallel with four to eight threads.

\subsection{In the browser}

The JavaScript decoder runs well. Interestingly, the \textbf{four-state}
variant is 7--13\% \emph{faster} than eight-state in JS, even on the same
hardware where eight states win in C. The reason is the JavaScript engine
(V8): it can keep four independent chains busy but not eight, so the extra
chains just add overhead without shortening the critical path. The eight-state
path still works in JS --- it lets one decoder accept C-encoded files without
conversion --- but four states is the sweet spot for the browser.

Concretely, a 9.35\,MB mammography image decodes in 19.8\,ms (473\,MB/s) using
eight workers with four-state strips. That means a web DICOM viewer can pull
a \texttt{.mic} file straight from object storage and decode it client-side,
with no server doing compression or format conversion. (These numbers are from
Node.js; full in-browser benchmarking is future work.)
Tables~\ref{tab:js-perf} and~\ref{tab:pics-js} give the JS numbers.

\begin{table}[htbp]
\centering
\caption{JavaScript decoder throughput on the ARM64 reference (Apple
M4~Pro, Node.js v22.16), for the four-state and eight-state FSE
variants. Median over 20 iterations after three warm-up runs.}
\label{tab:js-perf}
\begin{tabular}{lrrrrrr}
\toprule
 & & & \multicolumn{2}{c}{4-state} & \multicolumn{2}{c}{8-state} \\
\cmidrule(lr){4-5}\cmidrule(lr){6-7}
Image & Dimensions & Ratio & ms & MB/s & ms & MB/s \\
\midrule
MR  & $256{\times}256$   & 2.35$\times$ &   2.5 &  50 &   3.3 &  38 \\
CT  & $512{\times}512$   & 2.24$\times$ &  11.9 &  42 &  14.1 &  36 \\
CR  & $2140{\times}1760$ & 3.69$\times$ & 136.4 &  53 & 151.4 &  47 \\
MG1 & $2457{\times}1996$ & 8.79$\times$ &  68.1 & 137 &  78.3 & 120 \\
MG3 & $4774{\times}3064$ & 2.29$\times$ & 514.5 &  54 & 563.7 &  49 \\
\bottomrule
\end{tabular}
\end{table}

\begin{table}[htbp]
\centering
\caption{PICS parallel JavaScript decoder (\texttt{worker\_threads}) on
the ARM64 reference, with four-state and eight-state FSE per strip.
Worker count was swept over $\{1,2,4,8,14\}$; the best wall-time for
each variant is reported.}
\label{tab:pics-js}
\begin{tabular}{lrrrrr}
\toprule
 & & \multicolumn{2}{c}{4-state strips} & \multicolumn{2}{c}{8-state strips} \\
\cmidrule(lr){3-4}\cmidrule(lr){5-6}
Image & Strips &   ms & MB/s &   ms & MB/s \\
\midrule
MR  & 4 &   1.1 & 111 &   1.3 &  97 \\
CT  & 4 &   4.2 & 120 &   4.7 & 108 \\
CR  & 8 &  33.4 & 215 &  31.8 & 226 \\
MG1 & 8 &  19.8 & 473 &  20.5 & 457 \\
\bottomrule
\end{tabular}
\end{table}

% ============================================================
\section{Discussion}

\subsection{Reading the results together}

Three things tie the numbers together:

\begin{enumerate}
  \item \textbf{MIC's ratio edge over Zstandard} (5--22\%, 14\% average) comes
  mainly from working on 16-bit residuals directly instead of splitting them
  into bytes --- and that same choice removes the per-byte branches that slow
  decoding.
  \item \textbf{On the very largest mammography images on AMD64}, HTJ2K's block
  coder regains the speed lead by amortizing its setup over many pixels; MIC
  stays within 0--20\% there.
  \item \textbf{Multi-state decoding} is what makes MIC's decoder competitive
  with HTJ2K's vectorized block coder, and it costs nothing in file size.
\end{enumerate}

\subsection{Which codec to pick}

\begin{itemize}
  \item \textbf{Low-latency viewing (PACS, thumbnails, tomosynthesis
  playback):} single-threaded MIC is a sensible default (${\sim}600$--900\,MB/s
  per request); strip-parallel MIC for one big image at a time.
  \item \textbf{Smallest archival files:} JPEG-LS, which keeps ${\sim}10\%$
  better ratio --- at the cost of decoding 1.6--5.7$\times$ slower on ARM64.
  \item \textbf{DICOM transfer-syntax compatibility required:} HTJ2K, since
  it's a standard (and it's specifically strong on big mammography on AMD64).
  \item \textbf{No native code allowed (browsers, sandboxes):} the pure-Go
  reference or the 20\,KB JavaScript decoder, at about half the C decoder's
  speed.
\end{itemize}

A more structured comparison (a Pareto view of ratio vs.\ decode speed, plus a
weighted scoring matrix over ratio, speed, platform coverage, and deployment
ease) puts MIC and JPEG-LS on the efficiency frontier --- MIC for throughput,
JPEG-LS for ratio --- with HTJ2K and Zstandard dominated on those two axes.
HTJ2K still wins whenever standard compatibility is a hard requirement that
the scoring doesn't capture. Table~\ref{tab:decision-matrix} reports the
weighted matrix.

\begin{table}[htbp]
\centering
\caption{Codec selection matrix. Each criterion is normalised so the
best codec on that row scores~1.00. Quantitative rows are geometric
means across the 21-image set. Browser compatibility: shipped
purpose-built JS or WASM decoder used in production~=~1.00; WASM port of
the native library used in clinical viewers~=~0.85; no browser
path~=~0.0. Deployment ease: zero external runtime dependencies~=~1.00;
ubiquitous system dependency~=~0.90; C/C++ library requiring CMake build
and link integration~=~0.70. Weight columns $w_P$, $w_A$, $w_B$ reflect
typical priorities for (P)~PACS viewer / low-latency decode, (A)~archival
storage, (B)~browser delivery.}
\label{tab:decision-matrix}
\footnotesize
\setlength{\tabcolsep}{4pt}
\begin{tabular}{lrrrrrr}
\toprule
Criterion & $w_P$ & $w_A$ & $w_B$ & MIC & HTJ2K & JPEG-LS \\
\midrule
Compression ratio                & 0.20 & 0.50 & 0.10 & 0.91 & 0.90 & \textbf{1.00} \\
Decompression throughput (ARM64) & 0.30 & 0.10 & 0.15 & \textbf{1.00} & 0.53 & 0.31 \\
Encoding throughput (ARM64)      & 0.05 & 0.10 & 0.05 & \textbf{1.00} & 0.37 & 0.22 \\
AMD64 native                     & 0.10 & 0.05 & 0.05 & \textbf{1.00} & \textbf{1.00} & \textbf{1.00} \\
ARM64 native                     & 0.10 & 0.05 & 0.05 & \textbf{1.00} & \textbf{1.00} & \textbf{1.00} \\
Browser compatibility            & 0.05 & 0.00 & 0.40 & \textbf{1.00} & 0.85 & 0.85 \\
Deployment ease                  & 0.20 & 0.20 & 0.20 & \textbf{1.00} & 0.70 & 0.70 \\
\midrule
\textbf{Weighted score (P) viewer}   & --- & --- & --- & \textbf{0.98} & 0.74 & 0.69 \\
\textbf{Weighted score (A) archival} & --- & --- & --- & \textbf{0.96} & 0.78 & 0.80 \\
\textbf{Weighted score (B) browser}  & --- & --- & --- & \textbf{0.99} & 0.78 & 0.74 \\
\bottomrule
\end{tabular}
\end{table}

\subsection{Limitations}

\begin{itemize}
  \item \textbf{Small dataset.} 21 images across nine modalities is varied but
  far from a clinical archive; broader validation is needed.
  \item \textbf{Tuning on the test set.} The adaptive table-size thresholds and
  the predictor choice were informed by this dataset. The rules are simple (two
  thresholds, one predictor), which limits overfitting risk, but larger studies
  should confirm they generalize.
  \item \textbf{Baseline coverage.} Stronger byte-oriented baselines (Delta +
  byte/bit-shuffle + Zstd, DICOM RLE) would more sharply isolate the value of
  the 16-bit entropy stage; left to future work.
  \item \textbf{Reporting.} Tables give single representative throughput values
  rather than full confidence intervals.
  \item \textbf{Not yet a DICOM transfer syntax.} MIC works on the extracted
  pixel array; real PACS integration would need a private encapsulation or a
  standardization step.
\end{itemize}

% ============================================================
\section{Conclusion}

MIC is a 16-bit-native lossless codec for medical images: spatial prediction,
16-bit run-length coding, a large-alphabet table-based entropy coder, and a
multi-state decoder for speed. On 21 DICOM images it reaches an average
ratio of 3.46 (tied with HTJ2K, ${\sim}91\%$ of JPEG-LS) while being the
fastest decoder on 20/21 images on ARM64 and all 21 on AMD64, and the fastest
encoder on all 21 ARM64 images and 18/21 on AMD64. A 20\,KB JavaScript decoder
and a WebAssembly build show the format can be decoded client-side in a
browser.

\textbf{Future work:} larger multi-hospital datasets; stronger shuffle-based
baselines; full statistical reporting; an ARM64 vector kernel for the wavelet
variant; and benchmarking the JavaScript decoder in real browsers.

MIC is open source: \url{https://github.com/pappuks/medical-image-codec}

\end{document}
